\chapter{Understanding Satellite Image Bands and Tensors}
\label{chap:ch02}

\section{Multispectral Bands: Beyond Human Vision}
Human vision is trichromatic: our retinas contain photoreceptor cone cells sensitive to Red ($\approx 650\text{ nm}$), Green ($\approx 550\text{ nm}$), and Blue ($\approx 450\text{ nm}$) wavelengths. When viewing a digital photograph on a computer screen, each pixel is composed of these three primary additive channels (RGB).

Earth observation satellites, however, are not constrained by human biology. The Sentinel-2 Multispectral Instrument records radiation across 13 distinct bands, capturing crucial non-visible electromagnetic spectra:

\begin{table}[htbp]
\centering
\small
\begin{tabularx}{\textwidth}{llccX}
\toprule
\textbf{Band} & \textbf{Band Name} & \textbf{Center $\lambda$ (nm)} & \textbf{Resolution} & \textbf{Primary Physical Role} \\
\midrule
\texttt{B02} & Blue & 490 & 10 m & Atmosphere/aerosol scattering, clear water bodies \\
\texttt{B03} & Green & 560 & 10 m & Vegetation peak green reflectance, plant health \\
\texttt{B04} & Red & 665 & 10 m & Chlorophyll absorption in healthy vegetation \\
\texttt{B08} & Near-Infrared (NIR) & 842 & 10 m & Internal leaf cellular structure reflectance, water boundaries \\
\bottomrule
\end{tabularx}
\caption{The four 10-meter spatial resolution optical bands utilized in this project from Sentinel-2 Level-1C/2A products.}
\label{tab:s2_bands}
\end{table}

\subsection{Why the Near-Infrared (NIR) Band is Essential}
In the visible spectrum, healthy green leaves appear green because chlorophyll pigments absorb Red (\texttt{B04}) and Blue (\texttt{B02}) light to power photosynthesis, while reflecting a modest fraction of Green (\texttt{B03}) light ($\approx 10\text{--}15\%$).

In the Near-Infrared spectrum (\texttt{B08}), however, spongy mesophyll cells inside healthy leaves reflect up to $50\text{--}60\%$ of incoming radiation to prevent thermal overheating. Water bodies, on the other hand, absorb almost $100\%$ of NIR radiation, appearing jet black. Including \texttt{B08} allows our neural network to:
\begin{enumerate}
    \item Compute the Normalized Difference Vegetation Index: $\text{NDVI} = \frac{\text{B08} - \text{B04}}{\text{B08} + \text{B04}}$.
    \item Reliably distinguish vegetated crops from bare soil, and water coastlines from land, even when haze distorts visible colors.
\end{enumerate}

\section{From Photons to Digital Tensors}
When satellite sensors measure surface radiance, the values are initially recorded as integer \textbf{Digital Numbers} (DN). In Sentinel-2 processing pipelines:
\begin{equation}
\rho = \frac{\text{DN}}{10000.0}
\end{equation}
where $\rho \in [0, 1]$ represents dimensionless \textbf{Bottom-of-Atmosphere (BOA) Surface Reflectance}. A value of $0.0$ means total absorption, while $1.0$ represents total reflection.

\section{The Multimodal Input Tensor: \texorpdfstring{$[B, C, H, W]$}{[B, C, H, W]}}
In PyTorch and modern deep learning frameworks, multidimensional arrays of numbers are called \textbf{tensors}. Image batches are organized using the standard \textbf{channel-first} layout:
\begin{equation}
\mathbf{X} \in \mathbb{R}^{B \times C \times H \times W}
\end{equation}
Let us break down each dimension for our concrete project tensor: $[2, 7, 128, 128]$.

\begin{figure}[htbp]
\centering
\begin{tikzpicture}[
    font=\small\sffamily,
    band/.style={draw=black!70, thick, minimum width=3.4cm, minimum height=0.55cm, align=center, rounded corners=2pt},
    s1band/.style={band, fill=cyan!18},
    s2band/.style={band, fill=green!18},
    maskband/.style={band, fill=orange!22},
    outband/.style={band, fill=blue!15},
    arrow/.style={-Stealth, thick, draw=black!75}
]
% Left column: Modalities
\node[font=\bfseries\sffamily, align=center] at (0, 4.3) {7-Channel Stacked Input Tensor\\$\mathbf{X}_{\text{in}} \in \mathbb{R}^{B \times 7 \times 128 \times 128}$};

\node[s1band] (c0) at (0, 3.3) {Ch 0: S1 SAR VV $\in [0,1]$};
\node[s1band] (c1) at (0, 2.65) {Ch 1: S1 SAR VH $\in [0,1]$};
\node[s2band] (c2) at (0, 1.85) {Ch 2: S2 Red (B04) $\in [0,1]$};
\node[s2band] (c3) at (0, 1.20) {Ch 3: S2 Green (B03) $\in [0,1]$};
\node[s2band] (c4) at (0, 0.55) {Ch 4: S2 Blue (B02) $\in [0,1]$};
\node[s2band] (c5) at (0, -0.10) {Ch 5: S2 NIR (B08) $\in [0,1]$};
\node[maskband] (c6) at (0, -0.90) {Ch 6: Soft Mask $\mathbf{M} \in [0,1]$};

% Braces / Labels on left
\draw[decorate, decoration={brace, amplitude=5pt, mirror}, thick, cyan!60!black] (-1.9, 3.55) -- (-1.9, 2.4) node[midway, left=6pt, align=right, font=\scriptsize\sffamily] {\textbf{Branch 0 Input}\\$\mathbf{x}_{:, 0:2}$ ($2$ ch)};
\draw[decorate, decoration={brace, amplitude=5pt, mirror}, thick, green!50!black] (-1.9, 2.1) -- (-1.9, -1.15) node[midway, left=6pt, align=right, font=\scriptsize\sffamily] {\textbf{Branch 1 Input}\\$\mathbf{x}_{:, 2:7}$ ($5$ ch)};

% Decoupled Stem
\node[draw=cyan!70!black, fill=cyan!10, thick, rounded corners=4pt, minimum width=3.0cm, minimum height=1.1cm, align=center] (b0) at (4.5, 2.95) {\textbf{Stem Branch 0}\\\scriptsize\texttt{Conv2d(2, 16, k=7, p=3)}\\$[B, 16, 128, 128]$};
\node[draw=green!60!black, fill=green!10, thick, rounded corners=4pt, minimum width=3.0cm, minimum height=1.3cm, align=center] (b1) at (4.5, 0.55) {\textbf{Stem Branch 1}\\\scriptsize\texttt{Conv2d(5, 16, k=7, p=3)}\\$[B, 16, 128, 128]$};

\draw[arrow] (1.8, 2.95) -- (b0.west);
\draw[arrow] (1.8, 0.55) -- (b1.west);

% Concatenation & Output
\node[draw=purple!70!black, fill=purple!12, thick, rounded corners=4pt, minimum width=3.1cm, minimum height=1.4cm, align=center] (cat) at (8.6, 1.75) {\textbf{Channel Concat}\\\scriptsize\texttt{torch.cat(..., dim=1)}\\$\mathbf{F}_0 \in \mathbb{R}^{B \times 32 \times 128 \times 128}$\\\scriptsize (To Encoder Stage 0)};

\draw[arrow] (b0.east) -| (cat.north);
\draw[arrow] (b1.east) -| (cat.south);

\end{tikzpicture}
\caption{Channel layout and modality splitting inside \texttt{ECRformerAdaptedModel.forward()} and \texttt{DecoupledEncoder.forward()}. The 7-channel input tensor is sliced along dimension $C=1$ into 2 SAR channels (indices \texttt{0:2}) and 5 optical+mask channels (indices \texttt{2:7}), projected independently into 16 feature maps each, and concatenated into 32 channels.}
\label{fig:tensor_layout}
\end{figure}


\begin{itemize}
    \item \textbf{Batch Size ($B = 2$):} The number of independent geographical image patches processed concurrently in a single forward/backward pass on the GPU. Processing multiple patches simultaneously stabilizes gradient estimates.
    \item \textbf{Channels ($C = 7$):} The number of stacked physical feature maps:
    \begin{itemize}
        \item \texttt{Channel 0}: Sentinel-1 SAR VV backscatter, normalized to $[0, 1]$.
        \item \texttt{Channel 1}: Sentinel-1 SAR VH backscatter, normalized to $[0, 1]$.
        \item \texttt{Channel 2}: Sentinel-2 Red band (\texttt{B04}), reflectance in $[0, 1]$.
        \item \texttt{Channel 3}: Sentinel-2 Green band (\texttt{B03}), reflectance in $[0, 1]$.
        \item \texttt{Channel 4}: Sentinel-2 Blue band (\texttt{B02}), reflectance in $[0, 1]$.
        \item \texttt{Channel 5}: Sentinel-2 Near-Infrared band (\texttt{B08}), reflectance in $[0, 1]$.
        \item \texttt{Channel 6}: Cloud Mask $\mathbf{M} \in [0, 1]$, where $1.0$ denotes thick cloud and $0.0$ denotes clear sky.
    \end{itemize}
    \item \textbf{Height ($H = 128$) and Width ($W = 128$):} The spatial dimensions of each patch. At 10 meters per pixel, a $128 \times 128$ patch covers exactly $1.28\text{ km} \times 1.28\text{ km} = 1.6384\text{ km}^2$ of the Earth's surface!
\end{itemize}

\begin{mathbox}{Memory Footprint Calculation for a Batch}
How much VRAM does a single input batch consume?
For a single-precision 32-bit floating point tensor (\texttt{torch.float32}, 4 bytes per float):
\[
\text{Total elements} = B \times C \times H \times W = 2 \times 7 \times 128 \times 128 = 229{,}376\text{ floats}
\]
\[
\text{Memory (bytes)} = 229{,}376 \times 4\text{ bytes} = 917{,}504\text{ bytes} \approx 0.875\text{ MB}
\]
While the input itself is small ($<1\text{ MB}$), intermediate activations generated during transformer attention across multiple layers expand the GPU memory requirements during training to several gigabytes.
\end{mathbox}

\begin{codebox}{Tensor Slicing in \texttt{models/ecrformer\_adapted.py}}
In our verified implementation, the model receives this 7-channel tensor and splits it internally inside \texttt{ECRformerAdaptedModel.forward()}:
\begin{lstlisting}[language=Python]
# S1 SAR channels (Indices 0 and 1)
s1 = x[:, 0:2, :, :]          # Shape: [B, 2, H, W]

# Cloudy S2 optical channels (Indices 2, 3, 4, 5)
s2_cloudy = x[:, 2:6, :, :]   # Shape: [B, 4, H, W]

# Cloud mask channel (Index 6)
mask = x[:, 6:7, :, :]        # Shape: [B, 1, H, W]
\end{lstlisting}
Notice how clean Python array slicing preserves the separate physical modalities for specialized sub-networks!
\end{codebox}

\section{Summary of Chapter 2}
\begin{enumerate}
    \item Sentinel-2 provides multispectral data; our project utilizes Red (\texttt{B04}), Green (\texttt{B03}), Blue (\texttt{B02}), and Near-Infrared (\texttt{B08}) at 10 m resolution.
    \item Near-Infrared is vital for mapping vegetation vitality and sharp water boundaries that are hidden by clouds.
    \item Input data is formatted as a 4-dimensional tensor $[B, C, H, W]$ in channel-first format.
    \item The 7 channels consist of 2 SAR channels (VV, VH), 4 optical channels (RGB+NIR), and 1 cloud mask channel.
    \item Each $128 \times 128$ patch spans $1.64\text{ km}^2$ of terrain on Earth.
\end{enumerate}

\section{Practice Exercises}
\begin{enumerate}
    \item Calculate the total number of numbers in a tensor of shape $[4, 7, 256, 256]$.
    \item Why does a clear water lake appear dark in the Near-Infrared band (\texttt{B08}) but bright in the Blue band (\texttt{B02})?
    \item Write the PyTorch slicing syntax to extract only the Near-Infrared channel from an input tensor \texttt{x} of shape $[B, 7, H, W]$.
\end{enumerate}

\subsection*{Solutions}
\begin{enumerate}
    \item $4 \times 7 \times 256 \times 256 = 1{,}835{,}008$ numbers (approx $7.0\text{ MB}$ in float32).
    \item Liquid water absorbs long-wavelength near-infrared photons almost entirely (reflectance $\approx 0$), making lakes appear black in \texttt{B08}. Shorter blue wavelengths penetrate water or scatter off surface molecules, yielding measurable reflectance in \texttt{B02}.
    \item \texttt{nir = x[:, 5:6, :, :]} (or \texttt{x[:, 5, :, :]} to drop the channel dimension).
\end{enumerate}
